\subsection{A final classical evaluation-binding theorem}\label{sec:security14:soundness}

\begin{lemma}[Canonical bit opening is unique except under discrete-log failure]\label{lem:security14:canonical}
Fix the public Pedersen key $(g,h)$ and a collection
of tagged bit commitments. If two accepted
extracted witnesses assign different bits to
the same tagged group element, then from the
two witnesses one can compute $\alpha=\log_g h$.
If their bit values agree, their blinding
exponents agree as well.
\end{lemma}
\begin{proof}
For distinct bits, apply the explicit
collision formula of
Lemma~\ref{lem:realize:binding-reduction}.
For equal bits, cancel the shared
factor $g^b$ from the two
valid openings, giving
$h^r=h^{r'}$. Since $h$
has prime order $P$, the
blindings agree modulo $P$ by
Lemma~\ref{lem:global:unique-r}.
This is pointwise for every
tagged element, hence for the
whole registered bit vector.
\end{proof}

\begin{lemma}[A changed authenticated leaf gives a hash collision]\label{lem:security14:same-root}
Suppose the preprocessing extractor
obtains one valid indexed Merkle
commitment tree under root
$\mathsf{root}_j$, and a later
accepted relation proof or
authentication shell assigns
a different commitment vector
to an existing position under
the \emph{same} root.
Then a collision in the
underlying hash function can
be computed from the two
authenticated paths.
\end{lemma}
\begin{proof}
The preprocessing and later
witnesses provide the
different indexed leaf vectors
and their paths to the identical
root. The domain-separated
Merkle conflict reduction of
Lemma~\ref{lem:full:merkle-binding}
finds the first incompatible
serialized node inputs whose
hash outputs agree.
The differing inputs form
an explicit hash collision.
\end{proof}

\begin{lemma}[The preprocessing witness fixes one source table]\label{lem:security14:pre-fixed}
On the good event of the
prefix-preserving extraction
contract, the accepted
pre-evaluation transcript
determines one canonical
source table
$f_{\rm pre}:I_m\to\F$
before the point $z$ is chosen.
Every later accepted relation
that references the registered
source commitments uses the
same source table, unless a
discrete-log opening conflict
is obtained.
\end{lemma}
\begin{proof}
Definition~\ref{def:security14:extract-contract}
extracts actual canonical
bit openings for every source
field variable in the
preprocessing relation,
prior to revealing $z$.
Definition~\ref{def:full:field-encoding}
makes the source field value
uniquely determined by its
valid bit encoding.
Every subsequent use of that
source value refers to its
same named group commitments,
by Definition~\ref{def:full:registry}.
A different extracted source
value would change at least
one encoding bit and therefore
give two valid openings of
the same tagged group element.
Lemma~\ref{lem:security14:canonical}
then extracts the setup
discrete logarithm.
\end{proof}

\begin{theorem}[End-to-end classical evaluation binding: full-relation mode]\label{thm:security14:full-binding}
Under Definitions
\ref{def:security14:setup} and
\ref{def:security14:extract-contract},
for any classical PPT malicious
prover in the one-opening game,
the full-relation protocol of
Chapter~13 satisfies
\[
 \boxed{\Pr[\mathrm{accept}\ \wedge\
 v\neq\widetilde f_{\rm pre}(z)]
 \le\min\{1,\,
 \varepsilon_{\rm KE}
 +\varepsilon_{\rm DL}
 +\varepsilon_{\rm CRH}\}.}
\]
In particular, when the three
underlying advantages are negligible,
the protocol is classically
computationally evaluation binding.
There is no additional univariate
root-test or Reed--Solomon proximity
error in this full-relation
security bound.
\end{theorem}
\begin{proof}
Run the prescribed extractor
along the unchanged accepted
main path. Let $E$ be the
event of a failed extraction,
$D$ an extracted Pedersen
bit-opening conflict, and
$H$ an authenticated leaf
conflict for a fixed root.
Their probabilities are bounded
by the three displayed terms,
respectively. Outside
$E\cup D\cup H$, all
accepted proofs have valid
witnesses on one registered
source table $f_{\rm pre}$:
Lemma~\ref{lem:security14:pre-fixed}
fixes the source before $z$,
Lemma~\ref{lem:security14:canonical}
prevents conflicting
registered field values,
and Lemma~\ref{lem:security14:same-root}
prevents conflicting
authenticated leaf vectors.
The initial claimed-evaluation
relation in
Definition~\ref{def:full:protocol}
is an \emph{exact} circuit
equality for those same
registered source bits:
\[
 v=\sum_{a\in I_m}
 f_{\rm pre}(a)\operatorname{eq}_m(a,z).
\]
By the multilinear
interpolation identity
Lemma~\ref{lem:intro:interpolation},
the right side equals
$\widetilde f_{\rm pre}(z)$.
Thus an accepted false claim
cannot occur outside
$E\cup D\cup H$.
Apply the union bound and
cap by one. The proof never
replaces a complete-word
relation by a sparse-query
degree test; hence there is
no Schwartz--Zippel or
proximity term to add.
\end{proof}

\begin{corollary}[Knowledge-soundness of the expanded witness relation]\label{cor:security14:knowledge}
On the complement of the
extraction, discrete-log and
hash-collision events, an
accepted transcript determines
a single registered source
table $f_{\rm pre}$ and
witnesses for all accepted
NP relation statements that
verify the original masked
fold chain and terminal
Sumcheck relation.
\end{corollary}
\begin{proof}
Definition~\ref{def:security14:extract-contract}
returns a valid circuit
witness for every accepted
proof. The canonical commitment
and Merkle consistency
lemmas make their common
registered values agree.
The circuit relations in
Definitions
\ref{def:full:pre-relation} and
\ref{def:full:round-relations}
then give the fold chain,
degree budgets and masked
Sumcheck equalities.
The terminal relation is
verified by the final
extracted circuit witness.
All values refer to the
same source table by
Lemma~\ref{lem:security14:pre-fixed}.
\end{proof}

\begin{proposition}[Random challenges are unnecessary for full-relation algebraic binding]\label{prop:security14:uniform-not-needed}
Once the source has been
extracted before $z$ and
all complete relations are
enforced, the implication
\[
 \mathrm{accept}\Longrightarrow
 v=\widetilde f_{\rm pre}(z)
\]
outside the named cryptographic
bad events holds for every
fixed $z$ and every sequence
of legal fold challenges.
\end{proposition}
\begin{proof}
The proof of
Theorem~\ref{thm:security14:full-binding}
uses only the exact
registered-source identity
and multilinear
interpolation, not a
probabilistic root bound.
Both are valid for every
field input. Therefore
the conclusion holds
pointwise for any legal
challenge sequence.
\end{proof}

\begin{remark}[Why this is not yet an unconditional deployed theorem]
Theorem~\ref{thm:security14:full-binding}
is a \emph{rigorous reduction under an
explicit NP argument-of-knowledge
contract}. Chapter~13 provides the
circuits; Chapters~10--12 provide
a concrete classical affine component.
But an implemented proof system
for the \emph{full odd-field circuits}
with the prefix-preserving extraction
property has not yet been supplied
or machine audited. Assigning
$\varepsilon_{\rm KE}=0$ without
such an instantiation would be
unjustified.
\end{remark}
